\begin{landscape}
\footnotesize

\begin{longtable}{p{2cm} p{8cm} p{7cm} p{4cm} p{4cm}}

\caption{Gap Analysis Summary} \\

\toprule
\textbf{Gap} & \textbf{What existing methods do} & \textbf{What is missing} & \textbf{Implications for practice and theory} & \textbf{Supporting references} \\
\midrule
\endfirsthead

\multicolumn{5}{l}{\tablename\ \thetable\ -- Continued} \\
\toprule
\textbf{Gap} & \textbf{What existing methods do} & \textbf{What is missing} & \textbf{Implications for practice and theory} & \textbf{Supporting references} \\
\midrule
\endhead

\midrule
\multicolumn{5}{r}{Continued on next page} \\
\endfoot

\bottomrule
\endlastfoot

Gap 1: Absence of sentence-level ESG evidence extraction & Contemporary ESG NLP work often emphasizes document-level sentiment or coarse-grained topic signals (e.g., general tone of a sustainability report, or high-level topic modeling across entire documents) \cite{10.3390/su14159165,10.3390/su15054134,10.69593/ajbais.v4i04.130,10.1111/eufm.12509}. These approaches capture broad impressions but do not systematically attach sentiment to discrete ESG evidence sentences or clause-level textual cues. Some ABSA implementations target fine-grained aspect extraction within domains such as general text or non-financial domains, yet there is limited evidence of rigorously validated sentence-level ABSA pipelines specifically tuned to ESG disclosures, where regulatory-relevant evidence often resides in particular sentences, sentences within tables, or footnotes \cite{10.1002/csr.70089,10.3390/math12193119,10.2139/ssrn.3738629,10.1111/eufm.12509} Where sentence-level processing exists, it is frequently confined to generic sentiment at the sentence level or to shallow extraction of opinion-bearing phrases, without a formal linkage to precise ESG aspects, indicators, or regulatory lexicons. & A formal sentence-level ABSA framework for ESG filings that (i) maps each sentence to a concrete ESG aspect (e.g., emissions, governance oversight, supply-chain ethics), (ii) assigns sentence-level sentiment or evidentiary polarity with respect to that aspect, and (iii) propagates confidence metrics and explainability to regulators and investors. A standardized annotation schema and corpus for Indonesian and English ESG disclosures that anchors aspect terms to GRI/SASB/TCFD-aligned concepts at the sentence level, enabling reproducible model development and cross-study comparability. Evaluation protocols that quantify the precision, recall, and calibration of sentence-level aspect-sentiment associations, including robustness to multi-aspect sentences and to the presence of measured numerical indicators within the same sentence. & Without sentence-level evidence extraction, ABSA outputs remain insufficiently actionable for due diligence, risk assessment, and regulatory verification, because decision-makers require precise textual evidence linking sentiment signals to specific ESG topics and regulatory criteria \cite{10.1002/csr.70089,10.1108/mf-01-2023-0020,10.1111/eufm.12509} A sentence-level approach would enable more transparent risk signaling, facilitate traceability for greenwashing detection, and improve the interpretability of ABSA outputs through direct citation of textual evidence. & ABSA literature highlights the need for fine-grained, aspect-level extraction to reveal granular signals in ESG narratives \cite{10.1002/csr.70089,10.3390/math12193119,10.2139/ssrn.3738629,10.1111/eufm.12509} Cross-lingual ABSA and multilingual ESG contexts further motivate sentence-level analysis to capture language-specific expressions of risk or compliance indicators in Indonesian and English text \cite{10.1016/j.inffus.2025.103073,10.1111/eufm.12509,10.48550/arxiv.2511.00024} \\
\addlinespace
Gap 2: Conflation of tone and sentiment in existing approaches & Many studies treat ESG "tone" (the overall sentiment polarity of a document or section) as synonymous with "sentiment" toward ESG topics, effectively collapsing contextual cues into a single polarity score \cite{10.1002/csr.70089,10.69593/ajbais.v4i04.130,10.1108/mf-01-2023-0020} Some traditional sentiment analyses rely on bag-of-words or lexicon-based approaches that assign a global sentiment to documents, or to broad sections, without distinguishing whether sentiment reflects a positive environmental action, neutral governance reporting, or negative social impact within a specific topic \cite{10.3390/su14159165,10.1111/eufm.12509} ABSA literature emphasizes the distinction between sentiment polarity and opinion-bearing aspects, but practice in ESG contexts often lags in consistently separating topic-level sentiment from the broader document tone, particularly when sentences mention multiple ESG topics with conflicting cues. & A clear theoretical and methodological separation between tone (global document or section-level polarity) and targeted sentiment toward discrete ESG aspects, ensuring that ABSA outputs reflect aspect-specific sentiment rather than aggregate tone. Evaluation frameworks that can detect and quantify conflation errors, such as misattributing an English-positive sentence to an environmental aspect when the sentiment is actually about governance or labor practices. Multimodal and numerical signal integration to disambiguate tone from explicit sentiment about an ESG aspect, for example when numerical performance indicators accompany otherwise neutral text. & Conflation reduces the reliability of ABSA outputs for decision-makers who rely on precise signal attribution, such as due-diligence teams and regulators assessing material ESG risks and governance reforms. Clear separation improves model interpretability and supports explainable AI (XAI) approaches, enabling stakeholders to trace the rationale behind an aspect-level sentiment verdict. & Systematic ESG analytics reviews emphasize challenges around explainability, interpretability, and the need to map sentiments to concrete ESG dimensions \cite{10.1002/csr.70089,10.2139/ssrn.3738629,10.1108/mf-01-2023-0020} Cross-lingual ABSA research underscores the importance of maintaining topic-sentiment separation across languages to preserve semantic fidelity \cite{10.1016/j.inffus.2025.103073,10.1111/eufm.12509,10.48550/arxiv.2511.00024} \\
\addlinespace
Gap 3: LLM output instability across prompt templates & Large language models (LLMs) are increasingly applied to extract, summarize, and reason about ESG content via prompt-based or retrieval-augmented generation (RAG) paradigms. However, reported behavior often varies with prompt wording, template selection, and in-context examples, leading to instability in outputs across prompts or sessions \cite{10.1609/aaai.v38i21.30574,10.48550/arxiv.2511.14210} Some deployments adopt fixed prompts or single-shot prompts, which can yield inconsistent results when applied to diverse ESG documents, especially given domain-specific vocabulary and regulatory lexicons \cite{10.1609/aaai.v38i21.30574,10.48550/arxiv.2511.14210} There is growing recognition of the need for prompt engineering, prompt ensembles, and verification strategies to stabilize outputs and improve reproducibility \cite{10.1609/aaai.v38i21.30574,10.48550/arxiv.2511.14210} & Systematic studies that quantify prompt-induced variability for ESG ABSA tasks in Indonesian/English bilingual corpora, including metrics of stability (e.g., variance in aspect extraction, sentiment labeling, and regulatory mappings across templates). Robust prompt design guidelines and ensemble strategies (e.g., multiple prompts, self-consistency checks, majority voting) tailored to ESG ABSA with multilingual and legal/regulatory constraints. Integration of verification layers, such as post-hoc fact-checking against regulatory lexicons (GRI/SASB/TCFD) and cross-language alignment checks, to improve reliability of LLM outputs. & Without addressing prompt instability, LLM-based ABSA systems risk inconsistent outputs, undermining trust, regulatory audibility, and investor confidence. Stable, hybrid architectures combining LLMs with deterministic components (rule-based lexicons, supervised ABSA models) can reduce risk while preserving the benefits of generative capabilities. & The literature on LLMs in finance and ESG emphasizes model behavior, prompts, and the need for reliability and explainability, including considerations of bias, privacy, and regulatory compliance \cite{10.21203/rs.3.rs-6614220/v1,10.1142/9789811250903\_0004,10.30574/ijsra.2024.11.1.0305} Recent ESG document QA and RAG approaches illustrate how retrieval-augmented generation can improve accuracy and traceability, but also highlight instability risks in prompt-driven systems \cite{10.1609/aaai.v38i21.30574}. \\
\addlinespace
Gap 4: Neglect of OCR quality in document-level pipelines & Many ESG NLP studies assume textual input is available as clean machine-readable text, effectively bypassing the OCR bottleneck typical of scanned or PDF-disclosed documents. In practice, several document-analysis pipelines acknowledge OCR as a precursor step but do not systematically evaluate how OCR quality (character error rate, layout extraction accuracy) propagates into downstream ABSA performance, especially in multilingual settings where scripts and diacritics differ. Some studies of scanned or semi-structured documents from other domains (healthcare, historical archives) explicitly show that OCR quality, layout understanding, and document structure critically affect information extraction performance, suggesting similar implications for ESG reports that are often distributed as PDFs with embedded images or tables \cite{10.1109/icdar.2011.302,10.1093/jamiaopen/ooac045,10.21203/rs.3.rs-6644838/v1,10.3390/app122111063} & An end-to-end evaluation of how OCR quality, layout analysis, and document structure influence sentence-level ABSA and aspect-term extraction in Indonesian/English ESG disclosures. Techniques that couple OCR refinement (e.g., layout-aware OCR, region-based extraction, post-OCR text normalization) with ABSA pipelines, including strategies to handle multilingual content, tables, footnotes, and image-based figures that contain text. Standards for reporting OCR-related uncertainty in ESG analytics, including error-aware sentiment and evidence mapping to specific document sections. & Ignoring OCR quality leads to noise and bias in ABSA outputs, particularly in Indonesian contexts with mixed Indonesian/English text and domain-specific vocabulary. This undermines reliability for regulatory and investor use. A tightly integrated OCR-ABSA workflow with layout-aware processing improves signal fidelity and supports reproducibility, auditability, and scalability. & Document- and OCR-centric studies illustrate the impact of OCR and layout on downstream NLP tasks, including historical document digitization and healthcare document processing where layout information significantly affects extraction quality \cite{10.1108/jd-03-2023-0055,10.1093/jamiaopen/ooac045,10.1109/icdar.2011.302,10.21203/rs.3.rs-6644838/v1,10.1145/3151509.3151525,10.1145/3167132.3167236,10.1145/3322905.3322909} Multilingual and cross-language ABSA literature emphasizes the need for robust tokenization and alignment across languages, which are sensitive to OCR output quality in bilingual documents (Šmíd \& Král, 2025), (Gorovaia \& Makrominas, 2024), (Hang, 2025). \\
\addlinespace
Gap 5: Underdeveloped greenwashing detection methods & A subset of ESG NLP research focuses on identifying greenwashing signals by analyzing the volume, framing, and accessibility of environmental claims, sometimes using length, readability, copiousness of environmental content, or discrepancy between stated commitments and disclosed performance \cite{10.1111/eufm.12509,10.3389/frai.2025.1616007,10.3897/rio.6.e57602} Some studies explicitly link linguistic features to greenwashing risk, finding that longer or more favorable environmental narratives do not always correlate with substantive environmental performance, implying the need for evidence-based detection pipelines that tie textual claims to measurable indicators \cite{10.1111/eufm.12509} Cross-domain literature shows that detection of misleading sustainability content benefits from ABSA to locate sentiment around environmental topics and to cross-check with governance disclosures and quantified targets \cite{10.1111/eufm.12509,10.69593/ajbais.v4i04.130} & A rigorous, formally defined greenwashing detection framework that integrates sentence- or aspect-level ABSA with verifiable evidence from numerical performance data, regulatory disclosures, and third-party verification, specifically adapted to Indonesian and English ESG reporting. Methodologies that quantify inconsistency between qualitative statements and quantitative signals (e.g., emissions reductions, energy use intensity, board oversight metrics) and that produce explainable justifications for greenwashing flags. Benchmark datasets and evaluation protocols for greenwashing detection in bilingual ESG disclosures, including annotation guidelines that capture deceptive framing, selectivity in disclosures, or omissions. & Without robust greenwashing detection, investors and regulators face elevated risk of misinterpreting ESG commitments, potentially leading to misallocations of capital or regulatory blind spots. A multidisciplinary approach combining ABSA, factual verification, and explainable anomaly detection can strengthen trust in ESG communications and support more reliable sustainability governance. & Studies on greenwashing detection via NLP show that longer, more positive but less readable ESG content may indicate greenwashing risk, highlighting the need for deeper, evidence-aligned analyses beyond superficial sentiment \cite{10.1111/eufm.12509} ESG analytics literature emphasizes the importance of explainability and regulatory-aligned signals to counter greenwashing and enhance the credibility of ESG disclosures \cite{10.1002/csr.70089,10.2139/ssrn.3738629,10.1108/mf-01-2023-0020} \\


\end{longtable}

\end{landscape}